%% SYNTHETIC TEST UNIT generated by writing/tools/make_outline_book.py from FINAL_OUTLINE_v2 -- not book text
\chapter{Yizhuang Liu 的严格因子化研究}\label{ch:yzliu}
本单元为测试用的合成文本（由 FINAL\_OUTLINE\_v2 自动生成），陈述内容从略。

\section{第 1 节（测试）}\label{sec:yzliu-intro}

\begin{lemma}[部分指数化引理 \tagA]\label{lem:yzliu-partial-exp}
（测试用合成陈述：内容从略。）
\uses{prop:appB-webs, def:prelim-wilson}
\end{lemma}
\begin{proofsketch}测试。\end{proofsketch}
\paragraph{地位与缺口}测试：见文献。

\begin{proposition}[捏点奇点、超软模式与 $L\to\infty$ 极限 \tagA]\label{prop:yzliu-pinch-ultrasoft}
（测试用合成陈述：内容从略。）
\uses{lem:yzliu-partial-exp, def:lamet-quasi-tmd, def:lamet-soft-offlc}
\end{proposition}
\begin{proofsketch}测试。\end{proofsketch}
\paragraph{地位与缺口}测试：见文献。

\begin{assumption}[静态色单态 $Q\bar Q$ 对的谱假设 \tagC]\label{asm:yzliu-heavy-pair}
（测试用合成陈述：内容从略。）
地位：测试状态句。
文献：测试。
被使用于：\cref{cor:yzliu-soft-form-factor}。
解除途径：测试解除途径。
\uses{def:lamet-soft-offlc}
\end{assumption}

\begin{theorem}[离光锥软函数是快速色源对的形状因子 \tagA]\label{thm:yzliu-soft-hqet}
（测试用合成陈述：内容从略。）
\uses{def:lamet-soft-offlc, lem:yzliu-partial-exp, asm:lamet-aux-renorm}
\end{theorem}
\begin{proofsketch}测试。\end{proofsketch}
\paragraph{地位与缺口}测试：本陈述的地位由\cref{cor:yzliu-soft-form-factor}决定。

\section{第 2 节（测试）}\label{sec:yzliu-soft}

\begin{example}[单圈离光锥软函数与 $S_r$]\label{ex:yzliu-soft-oneloop}
（测试用合成陈述：内容从略。）
\uses{thm:yzliu-soft-hqet, lem:tmd-soft-oneloop, def:lamet-reduced-soft}
\end{example}

\begin{corollary}[软函数是静态色源对的形状因子（非微扰形式） \tagR]\label{cor:yzliu-soft-form-factor}
（$R_{\mathrm{cond}}$：在所列编号假设之下证明完整；结论的地位不高于其中最弱的假设。）
前提：\cref{asm:yzliu-heavy-pair}、\cref{thm:yzliu-soft-hqet}。
（测试用合成陈述：内容从略。）
非微扰地位：\cref{asm:intro-wightman}。
\uses{asm:yzliu-heavy-pair, thm:yzliu-soft-hqet}
\end{corollary}

\begin{definition}[双曲角与六种构型]\label{def:yzliu-hyperbolic-angle}
（测试用合成陈述：内容从略。）
\uses{def:lamet-soft-offlc, def:prelim-wilson}
\end{definition}

\begin{assumption}[维数正规化微扰论中的 Wightman 结构 \tagA]\label{asm:yzliu-wightman-pt}
（测试用合成陈述：内容从略。）
地位：测试状态句。
文献：测试。
被使用于：\cref{thm:yzliu-analyticity}、\cref{thm:yzliu-time-ordering}、\cref{conj:yzliu-complex-wilson}、\cref{thm:yzliu-jet-continuation}。
解除途径：测试解除途径。
\uses{def:prelim-qcd, def:prelim-brst}
\end{assumption}
\begin{proofsketch}测试。\end{proofsketch}

\section{第 3 节（测试）}\label{sec:yzliu-analyticity}

\begin{theorem}[软函数对复双曲角的解析性 \tagA]\label{thm:yzliu-analyticity}
（测试用合成陈述：内容从略。）
\uses{def:yzliu-hyperbolic-angle, asm:yzliu-wightman-pt, thm:yzliu-soft-hqet, def:tmd-soft}
\end{theorem}
\begin{proofsketch}测试。\end{proofsketch}
\paragraph{地位与缺口}测试：本陈述的地位由\cref{asm:yzliu-wightman-pt}决定。

\begin{proposition}[离光锥软函数中快度发散的因子化（线性性） \tagA]\label{prop:yzliu-rapidity-linear}
（测试用合成陈述：内容从略。）
\uses{def:lamet-soft-offlc, prop:tmd-evolution, lem:yzliu-partial-exp, prop:appB-webs}
\end{proposition}
\begin{proofsketch}测试。\end{proofsketch}
\paragraph{地位与缺口}测试：本陈述的地位由\cref{prop:front-rapidity-renorm}决定。

\begin{definition}[本征软函数的光锥定义]\label{def:yzliu-intrinsic-soft-lc}
（测试用合成陈述：内容从略。）
\uses{def:lamet-quasi-tmd, def:tmd-soft, rem:tmd-schemes}
\end{definition}

\begin{theorem}[本征软函数：定义的等价、方案无关性与反常维度 \tagA]\label{thm:yzliu-intrinsic-soft}
（测试用合成陈述：内容从略。）
\uses{def:yzliu-intrinsic-soft-lc, def:lamet-reduced-soft, prop:yzliu-rapidity-linear, thm:yzliu-analyticity, def:sudakov-cusp, prop:appB-cusp-oneloop}
\end{theorem}
\begin{proofsketch}测试。\end{proofsketch}
\paragraph{地位与缺口}测试：本陈述的地位由\cref{prop:yzliu-sI-twoloop}决定。

\section{第 4 节（测试）}\label{sec:yzliu-rapidity}

\begin{proposition}[$S_I$ 的两圈结构 \tagR]\label{prop:yzliu-sI-twoloop}
（$R_{\mathrm{PT}}$：在重整化微扰论中对红外正规化的部分子态严格。）
（测试用合成陈述：内容从略。）
\uses{def:yzliu-intrinsic-soft-lc, def:sudakov-cusp, prop:appB-cusp-oneloop}
\end{proposition}

\begin{example}[由 RGE 确定两圈 $L_b^2$ 系数]\label{ex:yzliu-sI-rge}
（测试用合成陈述：内容从略。）
\uses{prop:yzliu-sI-twoloop}
\end{example}

\begin{assumption}[光锥极限对 ${\rm Im}\,\theta$ 的一致性 \tagP]\label{asm:yzliu-lightcone-uniform}
（测试用合成陈述：内容从略。）
地位：测试状态句。
文献：测试。
被使用于：\cref{prop:yzliu-universality-classes}。
解除途径：测试解除途径。
\uses{def:yzliu-hyperbolic-angle, thm:yzliu-analyticity}
\end{assumption}

\begin{proposition}[软函数的普适类 \tagP]\label{prop:yzliu-universality-classes}
（测试用合成陈述：内容从略。）
\uses{thm:yzliu-analyticity, prop:yzliu-rapidity-linear, asm:yzliu-lightcone-uniform}
\end{proposition}
\paragraph{地位与缺口}测试：见文献。

\section{第 5 节（测试）}\label{sec:yzliu-time-ordering}

\begin{definition}[指数正规子下的软因子]\label{def:yzliu-exp-soft}
（测试用合成陈述：内容从略。）
\uses{def:tmd-soft, def:prelim-wilson}
\end{definition}

\begin{lemma}[坐标空间参数函数的正性（全子式矩阵树公式）]\label{lem:yzliu-matrix-tree}
（测试用合成陈述：内容从略。）
% uses: none
\end{lemma}

\begin{lemma}[末态求和恒等式]\label{lem:yzliu-final-state}
（测试用合成陈述：内容从略。）
% uses: none
\end{lemma}

\begin{theorem}[DY 形 TMD 软因子的时序无关性 \tagA]\label{thm:yzliu-time-ordering}
（测试用合成陈述：内容从略。）
\uses{def:yzliu-exp-soft, asm:yzliu-wightman-pt, lem:yzliu-matrix-tree, prop:appB-webs}
\end{theorem}
\begin{proofsketch}测试。\end{proofsketch}
\paragraph{地位与缺口}测试：本陈述的地位由\cref{asm:yzliu-wightman-pt}、\cref{conj:yzliu-complex-wilson}决定。

\section{第 6 节（测试）}\label{sec:yzliu-quasi-tmd}

\begin{proposition}[指数正规子下 DY 与 SIDIS 软因子相等 \tagP]\label{prop:yzliu-dy-sidis-exp}
（测试用合成陈述：内容从略。）
\uses{def:yzliu-exp-soft, lem:yzliu-final-state, thm:yzliu-time-ordering}
\end{proposition}
\paragraph{地位与缺口}测试：本陈述的地位由\cref{thm:yzliu-analyticity}、\cref{thm:tmd-sivers-sign}决定。

\begin{conjecture}[复时空中的 Wilson 圈 \tagC]\label{conj:yzliu-complex-wilson}
（测试用合成陈述：内容从略。）
\uses{asm:yzliu-wightman-pt, thm:yzliu-time-ordering}
\end{conjecture}
\paragraph{地位与缺口}测试：本陈述的地位由\cref{thm:yzliu-time-ordering}决定。

\begin{proposition}[准 TMD 因子化的 Collins 式证明：单线记号下的幂次计数 \tagA]\label{prop:yzliu-qtmd-powercount}
（测试用合成陈述：内容从略。）
\uses{def:lamet-quasi-tmd, thm:lamet-quasi-tmd, asm:lamet-staple, prop:yzliu-pinch-ultrasoft, thm:yzliu-intrinsic-soft, prop:tmd-regions}
\end{proposition}
\begin{proofsketch}测试。\end{proofsketch}
\paragraph{地位与缺口}测试：见文献。

\begin{proposition}[准 TMD 匹配核：RGE、重—轻分解与两圈结果 \tagA]\label{prop:yzliu-hard-kernel}
（测试用合成陈述：内容从略。）
\uses{thm:lamet-quasi-tmd, prop:yzliu-qtmd-powercount, def:sudakov-cusp, prop:lattice-renormalizability}
\end{proposition}
\begin{proofsketch}测试。\end{proofsketch}
\paragraph{地位与缺口}测试：见文献。

\section{第 7 节（测试）}\label{sec:yzliu-exclusive}

\begin{example}[单圈 $H=C^+_{HL}C^-_{HL}$ 与相位]\label{ex:yzliu-hl-decomposition}
（测试用合成陈述：内容从略。）
\uses{prop:yzliu-hard-kernel}
\end{example}

\begin{proposition}[准光前波函数振幅的因子化 \tagA]\label{prop:yzliu-lfwf}
（测试用合成陈述：内容从略。）
\uses{thm:lamet-quasi-tmd, prop:yzliu-hard-kernel, thm:yzliu-time-ordering}
\end{proposition}
\begin{proofsketch}测试。\end{proofsketch}
\paragraph{地位与缺口}测试：本陈述的地位由\cref{thm:yzliu-time-ordering}决定。

\begin{proposition}[大动量下的介子形状因子与 $S_I$ 的提取 \tagA]\label{prop:yzliu-meson-ff}
（测试用合成陈述：内容从略。）
\uses{prop:yzliu-lfwf, thm:yzliu-intrinsic-soft, lem:tmd-glauber}
\end{proposition}
\begin{proofsketch}测试。\end{proofsketch}
\paragraph{地位与缺口}测试：见文献。

\begin{proposition}[以准 TMD 表示的小 $q_T$ Drell–Yan 截面 \tagA]\label{prop:yzliu-dy-quasi}
（测试用合成陈述：内容从略。）
\uses{thm:tmd-factorization, thm:lamet-quasi-tmd, thm:yzliu-intrinsic-soft}
\end{proposition}
\begin{proofsketch}测试。\end{proofsketch}
\paragraph{地位与缺口}测试：本陈述的地位由\cref{thm:tmd-factorization}、\cref{thm:lamet-quasi-tmd}决定。

\section{第 8 节（测试）}\label{sec:yzliu-threshold}

\begin{proposition}[由准 TMD 之比得到 Sivers 不对称 \tagA]\label{prop:yzliu-sivers}
（测试用合成陈述：内容从略。）
\uses{thm:lamet-quasi-tmd, prop:lamet-tmd-ratios, thm:tmd-sivers-sign}
\end{proposition}
\begin{proofsketch}测试。\end{proofsketch}
\paragraph{地位与缺口}测试：见文献。

\begin{definition}[重—轻 Sudakov 因子与类空喷注函数]\label{def:yzliu-hl-jet}
（测试用合成陈述：内容从略。）
\uses{def:prelim-wilson, def:lattice-quasi, prop:yzliu-hard-kernel}
\end{definition}

\begin{theorem}[类空与类时重夸克喷注函数的解析延拓 \tagA]\label{thm:yzliu-jet-continuation}
（测试用合成陈述：内容从略。）
\uses{def:yzliu-hl-jet, asm:yzliu-wightman-pt}
\end{theorem}
\begin{proofsketch}测试。\end{proofsketch}
\paragraph{地位与缺口}测试：本陈述的地位由\cref{thm:yzliu-time-ordering}决定。

\begin{theorem}[准 PDF 匹配核的阈值因子化（Ji–Liu–Su） \tagA]\label{thm:yzliu-threshold}
（测试用合成陈述：内容从略。）
\uses{def:yzliu-hl-jet, thm:yzliu-jet-continuation, prop:yzliu-hard-kernel, prop:lamet-oneloop, thm:lattice-lamet}
\end{theorem}
\begin{proofsketch}测试。\end{proofsketch}
\paragraph{地位与缺口}测试：见文献。

\section{第 9 节（测试）}\label{sec:yzliu-renormalon}

\begin{example}[阈值因子化的单圈检验]\label{ex:yzliu-threshold-oneloop}
（测试用合成陈述：内容从略。）
\uses{thm:yzliu-threshold, prop:lamet-oneloop}
\end{example}

\begin{proposition}[反常维度代数 \tagA]\label{prop:yzliu-ad-algebra}
（测试用合成陈述：内容从略。）
\uses{thm:yzliu-threshold, prop:yzliu-hard-kernel, thm:yzliu-intrinsic-soft, thm:yzliu-jet-continuation}
\end{proposition}
\begin{proofsketch}测试。\end{proofsketch}
\paragraph{地位与缺口}测试：见文献。

\begin{proposition}[阈值因子化的 NNLO 检验 \tagP]\label{prop:yzliu-threshold-nnlo}
（测试用合成陈述：内容从略。）
\uses{thm:yzliu-threshold, prop:yzliu-hard-kernel, prop:lamet-twoloop}
\end{proposition}
\paragraph{地位与缺口}测试：见文献。

\begin{remark}[阈值重求和与大 $x$ 适用域]\label{rem:yzliu-threshold-domain}
（测试用合成陈述：内容从略。）
% uses: none
\end{remark}

\section{第 10 节（测试）}\label{sec:yzliu-gpd}

\begin{proposition}[重—轻 Sudakov 核的线性红外重整子 \tagP]\label{prop:yzliu-linear-renormalon}
（测试用合成陈述：内容从略。）
\uses{prop:yzliu-hard-kernel, def:power-borel, prop:power-renormalon-examples, prop:lamet-lrr}
\end{proposition}
\paragraph{地位与缺口}测试：本陈述的地位由\cref{prop:lattice-cs-kernel}决定。

\begin{conjecture}[线性重整子归一化的普适性 \tagC]\label{conj:yzliu-renormalon-norm}
（测试用合成陈述：内容从略。）
\uses{prop:yzliu-linear-renormalon}
\end{conjecture}
\paragraph{地位与缺口}测试：本陈述的地位由\cref{prop:yzliu-linear-renormalon}决定。

\begin{conjecture}[强子夸克双线性矩阵元的阈值渐近 \tagC]\label{conj:yzliu-threshold-hadronic}
（测试用合成陈述：内容从略。）
\uses{prop:yzliu-linear-renormalon, thm:yzliu-threshold}
\end{conjecture}
\paragraph{地位与缺口}测试：本陈述的地位由\cref{thm:yzliu-threshold}、\cref{asm:intro-bridge}决定。

\begin{proposition}[准 GPD 单圈匹配核在 $x=\pm\xi$ 附近的结构 \tagR]\label{prop:yzliu-gpd-oneloop}
（$R_{\mathrm{PT}}$：在重整化微扰论中对红外正规化的部分子态严格。）
（测试用合成陈述：内容从略。）
\uses{def:lamet-quasi-gpd, def:excl-gpd}
\end{proposition}

\section{第 11 节（测试）}\label{sec:yzliu-models}

\begin{conjecture}[GPD 在 $x=\pm\xi$ 附近的 LaMET 可计算性 \tagC]\label{conj:yzliu-gpd-xi}
（测试用合成陈述：内容从略。）
\uses{prop:yzliu-gpd-oneloop, thm:lamet-quasi-gpd}
\end{conjecture}
\paragraph{地位与缺口}测试：见文献。

\begin{theorem}[二维 $U(N)$ Gross–Neveu 模型中扭度展开的绝对收敛 \tagP]\label{thm:yzliu-gn-convergence}
（测试用合成陈述：内容从略。）
% uses: none
\end{theorem}
\paragraph{地位与缺口}测试：本陈述的地位由\cref{asm:lamet-uniform}决定。

\begin{proposition}[Gross–Neveu 模型中扭度展开的结构：重整子相消与真空凝聚 \tagP]\label{prop:yzliu-gn-structure}
（测试用合成陈述：内容从略。）
\uses{thm:yzliu-gn-convergence, def:power-borel}
\end{proposition}
\paragraph{地位与缺口}测试：见文献。

\begin{proposition}[超可重整二维 $O(N)$ 模型：动量空间 OPE 的阶乘增强 \tagP]\label{prop:yzliu-on-factorial}
（测试用合成陈述：内容从略。）
\uses{def:power-borel}
\end{proposition}
\paragraph{地位与缺口}测试：本陈述的地位由\cref{thm:lamet-hadronic}、\cref{rem:lamet-open}决定。

\section{第 12 节（测试）}\label{sec:yzliu-decay}

\begin{proposition}[二维非线性 $\sigma$ 模型：凝聚与跨级数 \tagP]\label{prop:yzliu-nlsm}
（测试用合成陈述：内容从略。）
\uses{def:power-borel}
\end{proposition}
\paragraph{地位与缺口}测试：本陈述的地位由\cref{prop:power-renormalon-examples}决定。

\begin{proposition}[二维大 $N_c$ QCD 中的准 PDF \tagP]\label{prop:yzliu-2dqcd}
（测试用合成陈述：内容从略。）
\uses{def:lattice-quasi, prop:power-thooft-model}
\end{proposition}
\paragraph{地位与缺口}测试：见文献。

\begin{remark}[坐标空间与动量空间之间变换的教训]\label{rem:yzliu-transforms}
（测试用合成陈述：内容从略。）
% uses: none
\end{remark}

\begin{remark}[瞬子与模型研究]\label{rem:yzliu-instantons}
（测试用合成陈述：内容从略。）
% uses: none
\end{remark}

\section{第 13 节（测试）}\label{sec:yzliu-status}

\begin{proposition}[标量场论中的谱表示与指数衰减（微扰论各阶） \tagP]\label{prop:yzliu-decay-extensions}
（测试用合成陈述：内容从略。）
% uses: none
\end{proposition}
\paragraph{地位与缺口}测试：本陈述的地位由\cref{thm:lamet-exp-decay}、\cref{lem:lamet-strip}、\cref{prop:lamet-truncation}决定。

\begin{remark}[与 CSS、Collins 与 SCET 的关系]\label{rem:yzliu-relation-css}
（测试用合成陈述：内容从略。）
% uses: none
\end{remark}

\begin{remark}[开放问题]\label{rem:yzliu-open}
（测试用合成陈述：内容从略。）
% uses: none
\end{remark}
